%%% FIELD proof-thm-gaussian-coordinate-lower
Put \(N=3n\).  By \(\textup{def:projected-experiment}\), concatenating the mutually independent folds \(D_K,D_D,D_A\) identifies \(Q_P\) with \(P^{\otimes N}\).  Thus every interval measurable with respect to the three folds is an interval in the \(N\)-observation experiment.

We first establish the nuisance-product lower bound on an index for which \(\varepsilon_{g,n}\varepsilon_{q,n}>0\).  Define
\[
 a_N=\frac{\varepsilon_{g,n}}{K},
 \qquad
 b_N=\frac{\varepsilon_{q,n}}{KA}.
\]
The constants \(A\) and \(K\) in \(\mathrm{CT2}\) do not depend on \(n\), so the product-smallness hypothesis in the theorem gives
\[
 a_Nb_N
 =\frac{\varepsilon_{g,n}\varepsilon_{q,n}}{K^2A}
 =o\!\left((\log N)^{-1/2}\right).
\]
Condition \(\mathrm{CT2}\) therefore permits application of \(\textup{lem:jms-support-localized-product-benchmark}\), with \(s=8\), source supplied fit \((\widehat g,\widehat q/A)\), and source radii \((a_N,b_N)\).  This invocation supplies only the nuisance-product/logarithmic mechanism from the two moment-matched fuzzy supports; no parametric term is attributed to that construction.

Let \(\widehat\theta\) be any real-valued estimator in the transformed \(N\)-observation experiment.  For source observations \(z_1,\ldots,z_N\), define
\[
 \widehat\vartheta_R(z_1,\ldots,z_N)
 =A^{-1}\widehat\theta(S_Az_1,\ldots,S_Az_N),
 \qquad S_A(x,t,y)=(x,t,Ay).
\]
Take \(\gamma=3/4\); since \(\alpha<1/2\), one has \(\gamma>\alpha\).  The cited support-localized benchmark yields \(i\in\{0,1\}\) and \(R\in\operatorname{supp}(\nu_{i,N})\) such that
\[
 R^{\otimes N}\!\left\{
 \left|\widehat\vartheta_R-\vartheta(R)\right|
 \ge c_\gamma a_Nb_N
 \left(\frac{\log(1/a_N)}{\log N}\right)^3
 \right\}\ge\gamma .
\]
Set \(P=R\circ S_A^{-1}\).  By \(\mathrm{CT2}\), this law satisfies
\[
 P\in\mathcal G_{\xi,n},
 \qquad
 (\widehat g,\widehat q)\in\mathcal H_n(P),
 \qquad
 \theta_0(P)=A\vartheta(R).
\]
Because \(S_A\) is a measurable bijection and \(Q_P=P^{\otimes N}\), the preceding probability is exactly
\[
 Q_P\!\left\{
 \left|\widehat\theta-\theta_0(P)\right|
 \ge c_\gamma A a_Nb_N
 \left(\frac{\log(1/a_N)}{\log N}\right)^3
 \right\}\ge\gamma .
\]
The use of the source radius \(b_N=\varepsilon_{q,n}/(KA)\) before the pushforward is essential, and it gives
\[
 A a_Nb_N
 =\frac{\varepsilon_{g,n}\varepsilon_{q,n}}{K^2}.
\]

Let \(\ell_n=1\vee\log(1/\varepsilon_{g,n})\).  On the present positive-product index, \(0<\varepsilon_{g,n}\le C_{\mathrm{env}}\).  Since \(K\ge e(1\vee C_{\mathrm{env}})\),
\[
 \log(1/a_N)=\log(K/\varepsilon_{g,n})\ge\ell_n.
\]
There is a universal constant \(c_0\ge1\) such that \(\log(3n)\le c_0\log(en)\) for every \(n\ge1\).  Consequently, using \(\textup{def:rate-scales}\),
\[
 A a_Nb_N
 \left(\frac{\log(1/a_N)}{\log N}\right)^3
 \ge c_1\varepsilon_{g,n}\varepsilon_{q,n}
 \min\!\left\{1,
 \left(\frac{\ell_n}{\log(en)}\right)^3\right\}
 =c_1\varepsilon_{g,n}\varepsilon_{q,n}\Lambda_n,
\]
where \(c_1>0\) is independent of \(n\).  Hence every estimator has error at least \(c_2\varepsilon_{g,n}\varepsilon_{q,n}\Lambda_n\), with probability at least \(\gamma\), at some law in the transferred fuzzy support.

By \(\mathrm{CT2}\), every law in that support belongs to \(\mathcal G_{\xi,n}\), is compatible with the supplied pair, and hence belongs to \(\mathcal U_n\) by \(\textup{def:gaussian-nondegenerate-subclass}\) and \(\textup{def:honesty-union}\).  The assumed union honesty of \(C_n\) therefore implies honesty over this transferred support.  Apply \(\textup{lem:honest-interval-length-reduction}\) with \(r=c_2\varepsilon_{g,n}\varepsilon_{q,n}\Lambda_n\), \(\beta=\gamma\), and this support as \(\mathscr R\).  Since \(\gamma-\alpha>0\), there is \(c_3>0\), independent of \(n\), such that
\[
 \sup_{P\in\mathcal G_{\xi,n}:(\widehat g,\widehat q)\in\mathcal H_n(P)}
 E_{Q_P}|C_n|
 \ge c_3\varepsilon_{g,n}\varepsilon_{q,n}\Lambda_n.
\]
If \(\varepsilon_{g,n}\varepsilon_{q,n}=0\), the same inequality is automatic because its right-hand side is zero; no fuzzy-support invocation is then needed.

We now obtain the parametric term by a separate compatible two-point experiment.  Set
\[
 \beta=\frac{\alpha+1/2}{2},
 \qquad \alpha<\beta<\frac12.
\]
Use the covariate law \(P_{X,n}\) and functions \(g_{*,n},q_{*,n}\) furnished by \(\mathrm{CT1}\).  For \(j\in\{0,1\}\), put
\[
 v_0=0,
 \qquad
 v_1=\Delta=\frac{h}{\sqrt N},
 \qquad
 P_j=P_{n,v_j},
\]
where \(h>0\) will be fixed independently of \(n\).  For all sufficiently large \(n\), \(\Delta\le h_0\), so \(\mathrm{CT1}\) gives
\[
 P_j\in\mathcal G_{\xi,n},
 \qquad
 (\widehat g,\widehat q)\in\mathcal H_n(P_j).
\]
Under \(P_j\), the defining equations in \(\mathrm{CT1}\) can be rewritten as
\[
 Y=v_jT+\{q_{*,n}(X)-v_jg_{*,n}(X)\}+A\zeta.
\]
Thus the partially linear coefficient is \(\theta_0(P_j)=v_j\), the structural nuisance is \(f_{0,j}=q_{*,n}-v_jg_{*,n}\), and the outcome disturbance is \(\xi_j=A\zeta\).  In particular, independence of \(\zeta\) from \((X,\eta)\), together with \(T=g_{*,n}(X)+\eta\), gives
\[
 E_{P_j}[\xi_j\mid X,T]=0,
 \qquad
 E_{P_j}[\xi_j^2\mid X,T]=A^2\ge c_\xi.
\]
The remaining membership conditions are already part of \(\mathrm{CT1}\).

Under \(P_0\), conditional division of the two Gaussian outcome densities gives the one-observation likelihood ratio
\[
 L=\frac{dP_1}{dP_0}
 =\exp\!\left\{
 \frac{\Delta\eta\{Y-q_{*,n}(X)\}}{A^2}
 -\frac{\Delta^2\eta^2}{2A^2}
 \right\}.
\]
Under \(P_0\), \(\{Y-q_{*,n}(X)\}/A=\zeta\) is standard Gaussian and independent of \(\eta\).  The standard-normal moment-generating identity therefore yields
\[
 E_{P_0}[L^2\mid X,\eta]
 =\exp(\Delta^2\eta^2/A^2).
\]
Since \(\eta\sim\mathcal N(0,1)\), direct Gaussian integration gives
\[
 E_{P_0}[L^2]
 =\left(1-\frac{2\Delta^2}{A^2}\right)^{-1/2}.
\]
Independence of the \(N\) observations now implies
\[
 \chi^2\!\left(P_1^{\otimes N},P_0^{\otimes N}\right)
 =\left(1-\frac{2h^2}{A^2N}\right)^{-N/2}-1.
\]
Choose \(h\le A/2\).  Then \(x=2h^2/(A^2N)\le1/2\), and the elementary inequality \(-\log(1-x)\le2x\), valid for \(0\le x\le1/2\), gives
\[
 \chi^2\!\left(P_1^{\otimes N},P_0^{\otimes N}\right)
 \le \exp(2h^2/A^2)-1.
\]
Reduce \(h>0\), if necessary, so that it also satisfies
\[
 \exp(2h^2/A^2)-1\le4(1-2\beta)^2.
\]
This choice depends only on \(\alpha\), \(A\), and the fixed feasibility radius \(h_0\).  By the definition of total variation and the Cauchy--Schwarz inequality,
\[
 \left\|P_1^{\otimes N}-P_0^{\otimes N}\right\|_{\mathrm{TV}}
 =\frac12E_{P_0^{\otimes N}}|L_N-1|
 \le\frac12\sqrt{\chi^2(P_1^{\otimes N},P_0^{\otimes N})}
 \le1-2\beta,
\]
where \(L_N=dP_1^{\otimes N}/dP_0^{\otimes N}\).

For an arbitrary estimator \(\widehat\theta\), define the disjoint events
\[
 A_j=\{|\widehat\theta-v_j|<\Delta/2\},
 \qquad j\in\{0,1\}.
\]
Their disjointness and the defining variational inequality imply
\[
 \begin{aligned}
 P_0^{\otimes N}(A_0)+P_1^{\otimes N}(A_1)
 &\le P_0^{\otimes N}(A_0)+P_0^{\otimes N}(A_1)
   +\left\|P_1^{\otimes N}-P_0^{\otimes N}\right\|_{\mathrm{TV}}\\
 &\le1+\left\|P_1^{\otimes N}-P_0^{\otimes N}\right\|_{\mathrm{TV}}.
 \end{aligned}
\]
It follows that
\[
 \max_{j\in\{0,1\}}
 P_j^{\otimes N}\!\left\{
 |\widehat\theta-v_j|\ge\frac{h}{2\sqrt N}
 \right\}
 \ge\frac{1-\|P_1^{\otimes N}-P_0^{\otimes N}\|_{\mathrm{TV}}}{2}
 \ge\beta.
\]
The two laws lie in the honesty union and are compatible with the supplied fits by \(\mathrm{CT1}\).  Applying \(\textup{lem:honest-interval-length-reduction}\) with \(\mathscr R=\{P_0,P_1\}\) and \(r=h/(2\sqrt N)\) yields a constant \(c_4>0\) such that
\[
 \sup_{P\in\mathcal G_{\xi,n}:(\widehat g,\widehat q)\in\mathcal H_n(P)}
 E_{Q_P}|C_n|
 \ge c_4N^{-1/2}
 \ge c_5n^{-1/2}.
\]

The same supremum is bounded below by the parametric term and by the nuisance-product/logarithmic term.  Using \(\max\{x,y\}\ge(x+y)/2\) and \(\textup{def:rate-scales}\), we conclude that
\[
 \sup_{P\in\mathcal G_{\xi,n}:(\widehat g,\widehat q)\in\mathcal H_n(P)}
 E_{Q_P}|C_n|
 \ge c\left\{n^{-1/2}
 +\varepsilon_{g,n}\varepsilon_{q,n}\Lambda_n\right\}
 =c d_n^-.
\]
Here \(c>0\) depends only on \(\alpha\), \(C_{\mathrm{env}}\), and \(c_\xi\), because \(A,K,h_0\) have precisely that dependence in \(\mathrm{CT1}\)--\(\mathrm{CT2}\).  Finally, \(\mathcal G_{\xi,n}\subseteq\mathcal G_n\) and \(\textup{def:length-vector}\) imply \(L_G(C_n)\ge c d_n^-\).  The argument uses the conditional outcome-variance restriction and \(\mathrm{CT1}\)--\(\mathrm{CT2}\) only to form these Gaussian lower-bound subexperiments; neither condition is asserted for the adaptive upper-frontier theorem.
